% !TEX root = ../algebra_02.tex

\section{Связь между алгебраической и геометрической кратностями собственного значения}

\begin{Theorem}{Связь геометрической и алгебраической кратностей}{}
	Для любого собственного значения $\lambda$ оператора $\alpha$ выполняется неравенство:
	\[
	1\le g_\lambda\le a_\lambda
	\]
	
\end{Theorem}

\begin{proof}
	Неравенство $g_\lambda\ge 1$ выполняется по определению собственного значения: так как существует хотя бы один ненулевой собственный вектор, размерность собственного подпространства $V_\lambda$ не меньше 1 (то есть $g_\lambda = \dim V_\lambda \ge 1$).
	
	Докажем второе неравенство. Пусть $g_\lambda = m$. Выберем базис $e_1, \dots, e_m$ в собственном подпространстве $V_\lambda$.
	Дополним его до базиса всего пространства $V$:
	\[
	E = (e_1, \dots, e_m, e_{m+1}, \dots, e_n).
	\]
	Поскольку векторы $e_1, \dots, e_m$ являются собственными для значения $\lambda$, действие оператора на них имеет вид $\alpha(e_i) = \lambda e_i$ для всех $i = 1, \dots, m$. 
	Следовательно, в базисе $E$ матрица оператора $\alpha$ имеет блочно-треугольный вид:
	\[
	[\alpha]_E = \begin{pmatrix} \lambda I_m & * \\ 0 & B \end{pmatrix},
	\]
	где $I_m$ — единичная матрица порядка $m$, а $B$ — некоторая квадратная матрица порядка $n-m$.
	
	Найдем характеристический многочлен оператора $\alpha$:
	\[
	\chi_\alpha(x) = \det([\alpha]_E - xI_n) = \det \begin{pmatrix} (\lambda - x)I_m & * \\ 0 & B - xI_{n-m} \end{pmatrix}.
	\]
	Определитель блочно-треугольной матрицы равен произведению определителей её диагональных блоков:
	\[
	\chi_\alpha(x) = \det((\lambda - x)I_m) \det(B - xI_{n-m}) = (\lambda - x)^m \det(B - xI_{n-m}).
	\]
	Таким образом, многочлен $(\lambda - x)^m$, то есть $(\lambda - x)^{g_\lambda}$, делит характеристический многочлен $\chi_\alpha(x)$. 
	Из определения алгебраической кратности как максимальной степени, в которой корень $\lambda$ входит в $\chi_\alpha(x)$, сразу следует, что $g_\lambda \le a_\lambda$.
\end{proof}